\clearpage
\label{sec:particelle}
\section{Dettaglio per particella}

Per ogni compresa e ogni particella vengono illustrate in dettaglio le caratteristiche del soprassuolo e il piano degli interventi silvicolturali.

Sono elencate in primo luogo le particelle governate a fustaia, corredate dei dati dendrometrici (numero di alberi per ettaro, volume per ettaro, area basimetrica per ettaro, altezza media, incremento percentuale di volume) suddivisi per specie e classe diametrica, nonché della ripresa attuale, anch'essa per specie.

Seguono le particelle governate a ceduo, per le quali è riportato il numero di ceppaie per ettaro.

\setlength{\LTleft}{0pt}
@@particelle(compresa=Serra,modello=particella-fustaia,ceduo=no)
@@particelle(compresa=Fabrizia,modello=particella-fustaia,ceduo=no)
@@particelle(compresa=Capistrano,modello=particella-fustaia,ceduo=no)

@@particelle(compresa=Serra,modello=particella-ceduo,ceduo=si)
@@particelle(compresa=Fabrizia,modello=particella-ceduo,ceduo=si)
@@particelle(compresa=Capistrano,modello=particella-ceduo,ceduo=si)
\clearpage
